\documentclass[10pt,a4paper]{amsart}
\usepackage[margin=19mm]{geometry}
\usepackage{amsmath,amssymb,mathtools}
\usepackage[scr=rsfs,cal=euler]{mathalfa}
\usepackage{xfrac}
\usepackage[unicode,hidelinks,bookmarks=false]{hyperref}
\newcommand{\ie}{i.e.\ }
\newcommand{\cf}{cf.\ }
\newcommand{\resp}{respectively\ }
\newcommand{\Subsection}{\subsection}
\providecommand{\nobreakdash}{\nobreak\hbox{-}}
\setlength{\emergencystretch}{2em}
\multlinegap0pt
\usepackage[shortlabels]{enumitem}
\usepackage{amssymb}
\newtheorem{theorem}{Theorem}
\newtheorem{lemma}{Lemma}
\title{\bf On a Lemma of Singh and Kumar}
\author{Jitender Singh}
\date{Source: arXiv:2505.08549v3 (CC BY 4.0)}
\begin{document}
\maketitle

\noindent\emph{Source and licence.} \bf On a Lemma of Singh and Kumar. Version arXiv:2505.08549v3 (CC BY 4.0), distributed by arXiv under the Creative Commons Attribution 4.0 International licence, http://creativecommons.org/licenses/by/4.0/, as recorded in the arXiv licence metadata for this exact version and shown on the abstract page https://arxiv.org/abs/2505.08549v3. Full licence notice: \texttt{licenses/arxiv-2505.08549-CC-BY-4.0.txt}.\par\smallskip

\noindent\textit{Editorial scope.} Counted unit: the article's theorem th:2 (source0.tex lines 137--145). The article's complete original proof of it is delivered with it (source0.tex lines 343--358, one \texttt{proof} environment of 16 lines, which the article places away from the statement). No mathematics is written, repaired or re-derived: the statement and the proof are the article's own lines, in the article's own notation, and no retained line is substituted or re-flowed.  Retained uncounted as the source-local context the proof needs: lines 102--110; lines 333--341.  Nothing was omitted from any retained span, so the package carries no omission record.  No retained line contains a citation, so the wrapper carries no bibliography and no bibliography entry.  No retained line contains a figure environment, an image-inclusion command or drawing code, so no image is delivered and the package declares no asset dependency.  Editorial formatting only, in addition: an AMS \texttt{amsart} preamble that restates the article's own declarations and macros used by the retained lines, namely the environments \texttt{theorem}, \texttt{lemma} on separate counters, together with the standard packages those lines need.  The retained environments print in this document as Theorem 1, Lemma 1 in order of appearance, whereas the article prints the same items with the numbering of its own full document.  The complete native source of the article as distributed in the arXiv e-print for this exact version is bundled unchanged as \texttt{sources/178049/source0.tex}.
\begin{theorem}\label{th:1}
Let $(R,v)$  be a discrete valuation domain. Let $f=a_0 + a_1x + \cdots + a_nx^n \in {R}[x]$ be a primitive polynomial such that there exist indices $j$ and $\ell$ with $1\leq \ell+1 \leq j\leq n$ for which the following hold.
\begin{enumerate}[label=$(\roman*)$]
\item $v(a_j)=0$,
\item $\frac{v(a_\ell)}{j-\ell}<\frac{v(a_i)}{j-i}$ for each $i=0,1,\ldots, j-1$ with $i\neq \ell$,
\item $\gcd(v(a_\ell),~j-\ell) = 1$.
\end{enumerate}
Then any factorization $f(x)=f_1(x)f_2(x)$ of $f$ in $R[x]$ has a factor of degree $\geq j-\ell$. In particular, if $j=n$ and $\ell=0$, then the polynomial $f$ is irreducible in $R[x]$.
\end{theorem}

\begin{lemma}\label{L:3}
Let $(R,v)$  be a discrete valuation domain. Let $f=a_0 + a_1x + \cdots + a_nx^n \in {R}[x]$ be such that there exists an index  $j$  with $1\leq j\leq n$ for which the following hold.
\begin{enumerate}[label=$(\roman*)$]
\item $v(a_j)=0$,
\item $\frac{v(a_0)}{j}<\frac{v(a_i)}{j-i}$ for each $i=1,\ldots, j-1$,
\item $\gcd(v(a_0),~j) = 1$.
\end{enumerate}
Then any factorization $f(x)=f_1(x)f_2(x)$ of $f$ in $R[x]$ satisfies $v(f_1(0))=v(a_0)$, or $v(f_2(0))=v(a_0)$.
\end{lemma}

\begin{theorem}\label{th:2}
Let $f=a_0 + a_1x + \cdots + a_nx^n \in \mathbb{Z}[x]$ be a primitive polynomial such that there exists a prime number $p$ and index $j$ with $1\leq j\leq n$ for which the following hold.
\begin{enumerate}[label=$(\roman*)$]
\item $v_p(a_j)=0$,
\item $\frac{v_p(a_0)}{j}<\frac{v_p(a_i)}{j-i}$ for each $i=1,\ldots,j-1$,
\item $\gcd(v_p(a_0),~j) = 1$.
\end{enumerate}
Let $d=a_0/p^{v_p(a_0)}$. If either $j=n$, or each zero $\theta$ of $f$ in ${\mathbb{C}}$ satisfies $|\theta|>d$, then  the polynomial $f$ is irreducible in $\mathbb{Z}[x]$.
\end{theorem}

\begin{proof}[\bf Proof of Theorem \ref{th:2}] The case $j=n$ follows from Theorem \ref{th:1} for $\ell=0$. Now assume that $j<n$. Suppose on the contrary that $f(x)=f_1(x)f_2(x)$ for nonconstant polynomials $f_1$ and $f_2$ in $\mathbb{Z}[x]$. Applying Lemma \ref{L:3} for $v=v_p$, the $p$-adic valuation of $\mathbb{Q}$, we have
 \begin{eqnarray}\label{e1}
v_p(f_1(0))=0,~\text{or}~v_p(f_2(0))=0,
 \end{eqnarray}
 which shows that $p$ does not divide at least one of $f_1(0)$ and $f_2(0)$. Assume that $p$ does not divide $|f_1(0)|$. If we define $v_p(a_0)=k$, then as $\pm p^kd=a_0=f(0)=f_1(0)f_2(0)$, it follows that $p^k$ divides $|f_2(0)|$, since $p\nmid d$. Consequently, we must have $|f_1(0)|\leq d$.

If $\theta_1,\ldots, \theta_{\deg f_1}$ are all zeros of $f_1$ in ${\mathbb{C}}$, then $|\theta_i|>d$ for all such $i$. If $0\neq \alpha \in \mathbb{Z}$ denotes the leading coefficient of $f_1$, then we may express $f_1$ as
 \begin{eqnarray*}
f_1(x)=\alpha(x-\theta_1)\cdots (x-\theta_{\deg f_1}).
\end{eqnarray*}
We then have
\begin{eqnarray*}
d\geq |f_1(0)|=|\alpha||\theta_1|\cdots |\theta_{\deg f_1}|>|\alpha_1| d^{\deg f_1}\geq d,
\end{eqnarray*}
which is a contradiction.
\end{proof}

\end{document}
